\documentclass[a4paper,11pt]{article}

% ----------------- Pacotes básicos -----------------
\usepackage[utf8]{inputenc}
\usepackage[T1]{fontenc}
\usepackage[french]{babel}
\usepackage{lmodern}
\usepackage{geometry}
\geometry{margin=2.5cm}
\usepackage{setspace}
\onehalfspacing

% ----------------- TikZ -----------------
\usepackage{tikz}
\usetikzlibrary{arrows.meta, positioning, shapes.geometric, fit}

% ----------------- Outros -----------------
\usepackage{hyperref}

\title{Projet Tetris multijoueur\\\large Analyse des cas d'utilisation}
\author{Équipe du projet}
\date{\today}

\begin{document}

\maketitle

\section{Introduction}

Ce document présente les principaux cas d'utilisation du système
\emph{Tetris multijoueur} développé dans le cadre du cours de
programmation orientée objet. L'objectif est de décrire, du point de
vue de l'utilisateur, les fonctionnalités offertes par le système et
les interactions avec les acteurs principaux.

\section{Acteurs}

Les acteurs identifiés sont les suivants :

\begin{itemize}
  \item \textbf{Joueur} : utilisateur humain du système. Il peut jouer
        en mode solo ou multijoueur, créer ou rejoindre des salles,
        et consulter les résultats d'une partie.
  \item \textbf{Serveur de parties} : composant logiciel chargé de
        gérer les salles multijoueurs, la synchronisation des
        parties et la communication entre les différents joueurs.
\end{itemize}

\section{Cas d'utilisation}

\begin{description}

  \item[UC1 -- Démarrer une partie solo]%
  \hfill\\
  \textbf{Acteur principal :} Joueur.\\
  Le joueur démarre une nouvelle partie de Tetris en mode solo. Le système
  initialise le plateau, la file de pièces et le score, puis passe à l'état
  \og partie en cours \fg.

  \item[UC2 -- Jouer en solo]%
  \hfill\\
  \textbf{Acteur principal :} Joueur.\\
  Le joueur contrôle les pièces (déplacement, rotation, chute rapide) pendant que
  le système fait tomber automatiquement la pièce courante, détecte les lignes
  complètes, met à jour le score et termine la partie en cas de \emph{game over}.

  \item[UC3 -- Créer une salle multijoueur]%
  \hfill\\
  \textbf{Acteurs principaux :} Joueur (hôte), Serveur de parties.\\
  Le joueur crée une nouvelle salle multijoueur (nom, nombre de joueurs, type de
  partie). Le serveur enregistre la salle et la rend visible pour les autres
  joueurs qui souhaitent la rejoindre.

  \item[UC4 -- Rejoindre une salle]%
  \hfill\\
  \textbf{Acteurs principaux :} Joueur, Serveur de parties.\\
  Le joueur consulte la liste des salles disponibles et demande à rejoindre une
  salle choisie. Le serveur vérifie les conditions (place disponible, partie
  non démarrée) et accepte ou refuse la demande.

  \item[UC5 -- Lancer la partie multijoueur]%
  \hfill\\
  \textbf{Acteurs principaux :} Joueur (hôte), Serveur de parties.\\
  L'hôte déclenche le début de la partie lorsque les joueurs requis sont
  présents. Le serveur notifie tous les clients et initialise un plateau pour
  chaque joueur, en passant la salle à l'état \og partie en cours \fg.

  \item[UC6 -- Jouer une partie multijoueur]%
  \hfill\\
  \textbf{Acteurs principaux :} Joueurs, Serveur de parties.\\
  Chaque joueur joue sa propre instance de Tetris. Lorsqu'un joueur efface des
  lignes, le client envoie un événement au serveur, qui calcule et distribue
  les effets (par exemple, lignes \og poubelle \fg) aux autres joueurs. La partie
  se termine lorsqu'il ne reste qu'un joueur actif ou lorsqu'une condition de fin
  prédéfinie est atteinte.

  \item[UC7 -- Voir les résultats de la partie]%
  \hfill\\
  \textbf{Acteur principal :} Joueur.\\
  À la fin d'une partie (solo ou multijoueur), le système affiche les
  résultats : score, nombre de lignes effacées, classement des joueurs, etc.
  Le joueur peut ensuite revenir au menu principal ou lancer une nouvelle
  partie.

\end{description}

\section{Diagramme de cas d'utilisation}

La figure~\ref{fig:usecase} présente une vue synthétique des cas
d'utilisation précédemment décrits.

\begin{figure}[h]
  \centering
  \begin{tikzpicture}[
    actor/.style   ={draw, rectangle, rounded corners, align=center,
                     minimum width=2.3cm, minimum height=0.9cm},
    usecase/.style ={draw, ellipse, align=center,
                     minimum width=3.8cm, minimum height=1.1cm},
    sysbox/.style  ={draw, rectangle, rounded corners,
                     inner sep=0.6cm}
  ]

    % ---------- Acteurs ----------
    \node[actor] (joueur)  at (-2,0) {Joueur};
    \node[actor] (serveur) at (10,0) {Serveur\\de parties};

    % ---------- Cas d'utilisation ----------
    % Mode solo (à gauche)
    \node[usecase] (startsolo) at (0,-2.5) {Démarrer une\\partie solo};
    % ligeiramente deslocado em x para que o include fique diagonal
    \node[usecase] (playsolo)  at (1.8,-4.7) {Jouer en solo};

    % Mode multijoueur (coluna à direita)
    \node[usecase] (createroom) at (7,-2.0) {Créer une salle\\multijoueur};
    \node[usecase] (joinroom)   at (7,-4.0) {Rejoindre une salle};
    \node[usecase] (startmp)    at (7,-6.0) {Lancer la partie\\multijoueur};
    \node[usecase] (playmp)     at (7,-8.0) {Jouer une partie\\multijoueur};
    \node[usecase] (results)    at (7,-10.0){Voir les résultats\\de la partie};

    % ---------- Boîte du système ----------
    \node[sysbox,
          fit=(startsolo) (playsolo) (createroom) (joinroom)
              (startmp) (playmp) (results),
          label=above:{Système : Tetris multijoueur}] (system) {};

    % ---------- Associations (sem sobreposição) ----------
    % Joueur
    \draw[-] (joueur.south) -- ++(0,-0.6) |- (startsolo.west);
    \draw[-] (joueur.south) -- ++(0,-0.6) |- (createroom.west);
    \draw[-] (joueur.south) -- ++(0,-0.6) |- (joinroom.west);
    \draw[-] (joueur.south) -- ++(0,-0.6) |- (results.west);

    % Serveur
    \draw[-] (serveur.south) -- ++(0,-0.6) |- (createroom.east);
    \draw[-] (serveur.south) -- ++(0,-0.6) |- (joinroom.east);
    \draw[-] (serveur.south) -- ++(0,-0.6) |- (startmp.east);
    \draw[-] (serveur.south) -- ++(0,-0.6) |- (playmp.east);

    % ---------- Relações <<include>> (não sobrepostas) ----------
    % diagonal leve entre os nós para o texto não ficar em pé
    \draw[dashed, -{Stealth[length=4pt]}]
      (startsolo.south east) ..
      controls +(0.5,-0.2) and +(-0.5,0.2) ..
      node[above,sloped]{\scriptsize <<include>>}
      (playsolo.north west);

    \draw[dashed, -{Stealth[length=4pt]}]
      (startmp.south) ..
      controls +(0,-0.7) and +(0,0.7) ..
      node[right]{\scriptsize <<include>>}
      (playmp.north);

    \draw[dashed, -{Stealth[length=4pt]}]
      (playmp.south) ..
      controls +(0,-0.7) and +(0,0.7) ..
      node[right]{\scriptsize <<include>>}
      (results.north);

  \end{tikzpicture}
  \caption{Diagramme de cas d'utilisation du système Tetris multijoueur.}
  \label{fig:usecase}
\end{figure}

\end{document}
